% V2扫描PDF141--142原编号思考题17.1--17.12；需要图17.31--17.37独立审核。
\section*{第17章\quad 磁场和它的源：思考题}
\paragraph{思考题 17.1}
一根电源引线和对应返回线为什么通常应拧在一起以降低对电子仪器的磁干扰？

\paragraph{思考题 17.2}
两根互相绝缘的相交直导线通过大小相等的电流，原图17.31所示两支成直角，在什么点合磁场为零？

\paragraph{思考题 17.3}
一根导线分为菱形电路上下两条对称分支，原图17.32中菱形两对角线交点处磁场方向与大小如何？

\paragraph{思考题 17.4}
解释等离子体电流产生自磁场，受到径向洛伦兹力后发生向轴心收缩的箍缩效应（Z-pinch）。

\paragraph{思考题 17.5}
在托卡马克中，螺绕环线圈和等离子体本身的感应环流怎样共同形成螺旋磁力线，对等离子体有什么约束作用？

\paragraph{思考题 17.6}
任取包围条形磁铁一个磁极的闭合曲面，穿过它的磁通量是否会不为零？

\paragraph{思考题 17.7}
磁场是否为保守场？区分 (
ablacdotoldsymbol B=0) 与 (
abla	imesoldsymbol B
e0)。

\paragraph{思考题 17.8}
在无传导电流的空间区域，若磁力线是平行直线，磁场大小是否必然处处相同？说明静磁真空条件。

\paragraph{思考题 17.9}
证明静磁条件下两磁极间的磁感应线不能按原图17.35那样突然终止于空气中。

\paragraph{思考题 17.10}
原图17.36一根长密绕螺线管通过电流I，闭合路径L局部绕过管体，求 (oint_Loldsymbol Bcdot doldsymbol l) 对穿链电流的依赖。

\paragraph{思考题 17.11}
对原图17.37不规则截面的长密绕螺线管，内部磁场是否一定均匀？是否可以无条件套用 (B=mu_0 nI)？

\paragraph{思考题 17.12}
参照原图17.29充电、放电电容器，电容板间位移电流及周围环绕磁场的方向分别如何变化？
